\documentclass[aspectratio=169]{beamer}

\mode<presentation>
{
  \usetheme{default}
  \usecolortheme{default}
  \usefonttheme{default}
  \setbeamertemplate{navigation symbols}{}
  \setbeamertemplate{caption}[numbered]
  \setbeamertemplate{footline}[frame number]  % or "page number"
  \setbeamercolor{frametitle}{fg=white}
  \setbeamercolor{footline}{fg=black}
}

\usepackage[english]{babel}
\usepackage[utf8x]{inputenc}
\usepackage{tikz}
\usepackage{courier}
\usepackage{array}
\usepackage{bold-extra}
\usepackage{minted}
\usepackage[thicklines]{cancel}
\usepackage{fancyvrb}
\graphicspath{{figures/}}

\xdefinecolor{dianablue}{rgb}{0.18,0.24,0.31}
\xdefinecolor{darkblue}{rgb}{0.1,0.1,0.7}
\xdefinecolor{darkgreen}{rgb}{0,0.5,0}
\xdefinecolor{darkgrey}{rgb}{0.35,0.35,0.35}
\xdefinecolor{darkorange}{rgb}{0.8,0.5,0}
\xdefinecolor{darkred}{rgb}{0.7,0,0}
\definecolor{darkgreen}{rgb}{0,0.6,0}
\definecolor{mauve}{rgb}{0.58,0,0.82}
\xdefinecolor{lightyellow}{rgb}{1.0,1.0,0.75}

\title[June 29th, 2022]{%
Modern Python analysis ecosystem for High Energy Physics%
}
\author{Jim Pivarski, Matthew Feickert, Gordon Watts}
\institute{Princeton University, University of Wisconsin-Madison, University of Washington}
\date[June 29th, 2022]{The Python Exchange for DOE Employees\\June 29th, 2022}

\usetikzlibrary{shapes.callouts}

\begin{document}

\logo{\pgfputat{\pgfxy(0.11, 7.4)}{\pgfbox[right,base]{\tikz{\filldraw[fill=dianablue, draw=none] (0 cm, 0 cm) rectangle (50 cm, 1 cm);}\mbox{\hspace{-8 cm}\raisebox{0.1 cm}{\includegraphics[height=0.8 cm]{iris-hep-logo-long.pdf}}\hspace{0.1 cm}}}}}

\begin{frame}
  \titlepage
\end{frame}

\logo{\pgfputat{\pgfxy(0.11, 7.4)}{\pgfbox[right,base]{\tikz{\filldraw[fill=dianablue, draw=none] (0 cm, 0 cm) rectangle (50 cm, 1 cm);}\mbox{\hspace{-8 cm}}}}}

% Uncomment these lines for an automatically generated outline.
%\begin{frame}{Outline}
%  \tableofcontents
%\end{frame}

% START

\include{src/matthew.tex}
%
\include{src/jim.tex}
%
\include{src/gordon.tex}

\begin{frame}{Conclusions: Modern Pythonic analysis in HEP is happening now}
\Large
\begin{center}
A confluence of scientific tools and scientists has lead to a feature-complete \textbf{Scikit-HEP} in the last 5 years
\end{center}
\vspace{0.1 cm}

\begin{center}
\href{https://indico.cern.ch/event/613842/contributions/2591057/}{\includegraphics[width=0.625\linewidth]{scikit-hep-in-2017.png}}
\end{center}

\vspace{-1.5 cm}
\hfill \mbox{Eduardo Rodrigues\hspace{-0.5 cm}}
\vspace{1.5 cm}
\end{frame}

\begin{frame}{Conclusions: Modern Pythonic analysis in HEP is happening now}
\Large
\begin{center}
Growing PyHEP \textbf{ecosystem} is enabling analysts in HEP to explore and reduce the time to insight
\end{center}
\vspace{0.1 cm}

\begin{center}
  \href{https://indico.cern.ch/event/1140031/}{\includegraphics[width=0.55\linewidth]{shells-hep.pdf}}
\end{center}
\end{frame}

\end{document}
